%======================= PREAMBOLO DICHIARAZIONI INIZIALI===============%
\documentclass[10pt,oneside,a4paper]{article}

\usepackage[latin1]{inputenc} 
\usepackage[italian]{babel}
\usepackage{siunitx} %Inserisce automaticamente i dati con le unità  di misura correttamente formattate del SI (utilizzo: \SI{0.82}{m^2}, in generale \SI{misura con il punto decimale}{unità  di misura})
\sisetup{output-decimal-marker = {.}, separate-uncertainty = true, input-uncertainty-signs = \pm, detect-weight=true, detect-family=true} %per usare SI con il punto decimale
\usepackage{listings} %Per citare codice informatico formattandolo correttamente
\usepackage{amsmath}
\usepackage{graphicx}
\usepackage{geometry}
\usepackage{epigraph}
\usepackage{booktabs}	%tabelle migliorate
\usepackage{tablefootnote}	%note a piè di pagina in tabella
\usepackage{threeparttable} %tabella con note a piè di tabella
\usepackage{caption}	%descrizione per figure
\captionsetup{tableposition=top,figureposition=bottom,font=small} %setup descrizione
\usepackage{float}
\usepackage{esvect} %vettori
\usepackage{longtable} %tabelle lunghe
\usepackage[dvipsnames]{xcolor}
\definecolor{sepia}{HTML}{80002A}
\usepackage[colorlinks=true, citecolor=black, linkcolor=sepia, urlcolor=black]{hyperref}
\usepackage{mathrsfs}
\usepackage[utf8]{inputenc}


\usepackage{listings} %Per inserire codice
\lstnewenvironment{codice_c}[1][]
{\lstset{basicstyle=\small\ttfamily, columns=fullflexible,
keywordstyle=\color{red}\bfseries, commentstyle=\color{blue},
language=C, basicstyle=\small,
numbers=left, numberstyle=\tiny,
stepnumber=2, numbersep=5pt, frame=shadowbox, #1}}{}

\setcounter{section}{-1}

%========= PRIMA PAGINA ===========%
\title{\textsc{Misura della frequenza di fenomeni radioattivi e studio di un pallinometro}}
\author{\small{G. Galbato Muscio} \and \small{L. Gravina} \and \small{L. Graziotto} \and \small{M. Rescigno}}
\date{}

\begin{document}
	\begin{figure}
		\centering
		\includegraphics[scale=0.5, trim={2.8cm 8.9cm 0 9cm}, clip]{logo.png}
	\end{figure}
	\maketitle
	\begin{center} 
		\fbox{{\fontsize{12pt}{8mm}\textsc{Gruppo B2.3}}} \\
		\vspace{1cm}
		\begin{tabular}{ccc}
			Esperienza di laboratorio && Consegna della relazione \\
			\emph{\small{3 maggio 2017}} && \emph{\small{14 maggio 2017}} \\
										&& 	\emph{\small{23 maggio 2017} (seconda parte)}  \\
		\end{tabular} 
		
		\vspace{0.5cm}
		
	\end{center}
\hrule
\vspace{0.5cm}
\begin{abstract}
Utilizzando un contatore Geiger-M\"uller misuriamo la radioattività  ambientale, quella dovuta ai raggi cosmici e la radiazione emessa da un blocco di tufo. Studiamo quindi la distribuzione delle biglie in un pallinometro fisico e in uno simulato.	
\end{abstract}
\newpage
\tableofcontents %Indice
\listoftables %Indice delle tabelle
\listoffigures %Indice dei grafici
\pagebreak
\section{Convenzioni e formule}
In questa relazione verranno usate le seguenti convenzioni:
\begin{enumerate}
	\item sarà usato il punto [ $.$ ] come separatore decimale;
	\item l'approssimazione decimale della cifra $5$ sarà fatta per eccesso;
	\item al fine di snellire la relazione e migliorarne la leggibilità , riporteremo nel corpo del documento solamente le tabelle riepilogative e alcuni grafici, e dedicheremo un'appendice finale alle tabelle e ai grafici più dettagliati.
\end{enumerate}
Inoltre, si farà riferimento alle seguenti formule:
\begin{enumerate}
	\item media 
	\begin{equation}\label{eq:media}
	\bar{x} = \frac{1}{N}\sum_{i=1}^Nx_i;
	\end{equation}
	\item varianza
	\begin{equation}\label{eq:varianza}
	\sigma^2 = \frac{1}{N}\sum_{i=1}^N(x_i-\bar{x})^2;
	\end{equation}
	\item deviazione standard
	\begin{equation}\label{eq:deviazione}
	\sigma = \sqrt{\sigma^2}.
	\end{equation}	
\end{enumerate}


%===============SCOPO E DESCRIZIONE DELL'ESPERIENZA==============%
\section{Scopo e descrizione dell'esperienza}
\label{sec:descrizione}



%================APPARATO SPERIMENTALE======================%		
\section{Apparato Sperimentale}


%==================SEQUENZA OPERAZIONI SPERIMENTALI==========
\section{Sequenza Operazioni Sperimentali} 

\subsecion{Estrazione del momento della forza d'attrito e del momento d'inerzia}

\subsecion{Studio della perdita di energia potenziale gravitazionale a causa dell'attrito}

\subsection{studio del moto del volano in presenza di attrito del mezzo}

%========================= CONCLUSIONI =============
\section{considerazioni Finali}

\pagebreak
%=================APPENDICE==================================

\section{Appendice: tabelle e grafici}


\end{document}
